\documentclass[11pt]{article}

\usepackage[a4paper,margin=25mm]{geometry}
\usepackage{amsmath,amssymb}
\usepackage{booktabs,longtable,array}
\usepackage{caption}
\usepackage[hidelinks]{hyperref}
\usepackage[T1]{fontenc}
\usepackage{microtype}
% Long identifiers such as \texttt{_accumulate_gene_split_event_vjp_kernel} do
% not fit on one line; [htt] lets them hyphenate inside typewriter text.
\usepackage[htt]{hyphenat}
\usepackage{listings}

\setlength{\emergencystretch}{3em}
\captionsetup[table]{skip=4pt,font=small}

\lstset{
  basicstyle=\ttfamily\small,
  breaklines=true,
  frame=single,
  framesep=4pt,
  columns=fullflexible,
  keepspaces=true,
}

% A source location: file and line, rendered without hyphenation.
\newcommand{\loc}[1]{\nolinkurl{#1}}
% A directory or path.
\newcommand{\dir}[1]{\nolinkurl{#1}}

\title{Code Review of the \texttt{main} Branch of \texttt{gpurec}}
\author{Review carried out with four parallel review agents}
\date{11 August 2026}

\begin{document}
\maketitle

\begin{abstract}
This is a review of the \texttt{main} branch of \texttt{gpurec} as of commit
\texttt{358a5b80}.  It answers six questions: how much code there is, whether
any files are bloated, whether the code is clean, whether the documentation is
good, whether the naming conventions are consistent, and whether the directory
tree is well organised.  Every count and every file-and-line reference was
produced by reading or running the code.  Nothing in the repository was
modified.  Section~\ref{sec:summary} is a summary and a list of what to fix, in
order; a reader who reads only that page will have the actionable content.
\end{abstract}

\tableofcontents

%======================================================================
\section{Summary and what to fix}
\label{sec:summary}

\subsection*{Verdict}

\textbf{This branch is not yet fit to be handed to an outside user, and the
reason is not the code --- it is everything around the code.}  The discipline
inside the source files is unusually good: no leftover markers, no
commented-out blocks, no stray debug output, no unsafe Rust, and naming
conventions that hold across 358 functions.  What fails is the surface a new
user touches.  The one shipped example does not run, the rule needed to make it
run is written down nowhere, there are no tests, and sixteen places in the
source point at documentation files that were deleted.

For reading and for small local fixes, this is a comfortable codebase.  For any
change that moves a number, it is not, because there is nothing on this branch
that can tell you whether you changed an answer.

\subsection*{The five findings that matter}

\begin{enumerate}
  \item \textbf{The shipped example does not run, and the rule it breaks is
        undocumented.}  A gene leaf must be named \texttt{SPECIES\_gene}; the
        code splits on the first underscore.  That rule appears in no README,
        docstring or comment.  The shipped example uses leaves called
        \texttt{a} and \texttt{b}, so the headline snippet in the README fails.
        A second file, \dir{examples/tiny/gene.map}, is presented by two
        READMEs as the thing that maps genes to species and is read by nothing.
        (Section~\ref{sec:docs-example}.)
  \item \textbf{There are no tests and no automated checks, and three files
        claim otherwise.}  The most damaging is a comment naming a ``permanent
        guard test'' that would catch the settings file drifting out of step
        with the Python code it mirrors.  That test is not on this branch, so
        the two copies are free to disagree silently.
        (Section~\ref{sec:no-tests}.)
  \item \textbf{One stale default is wrong by a factor of 21.}  A function
        declares \texttt{neumann\_terms=3} where the configured value is 64.
        Production callers always pass it explicitly, so it never fires in
        normal use --- it fires for whoever calls that function by hand.
        (Section~\ref{sec:stale-default}.)
  \item \textbf{The same settings live in three or four places, and six files
        take settings from shell environment variables.}  Two of those
        environment variables are read \emph{before} the argument the caller
        explicitly passed, so a deliberate choice in code is silently
        overridden by something invisible.  (Sections~\ref{sec:config-mirrors}
        and~\ref{sec:envvars}.)
  \item \textbf{986 lines have no caller, and two settings classes sit in the
        wrong package.}  The misplaced settings classes are why the settings
        package imports the model package, why the lowest-level kernels import
        the model package, and why one import had to be deferred inside a
        function with a nineteen-line comment apologising for it.
        (Sections~\ref{sec:dead-code} and~\ref{sec:imports}.)
\end{enumerate}

\subsection*{What to fix, in order}

Effort is a rough order of magnitude for one engineer who knows the code.

\begin{center}
\captionof{table}{Recommended order of work.  ``Owner'' marks items that need a
decision from whoever owns the branch rather than a code change.}
\begin{longtable}{@{}rp{0.60\textwidth}ll@{}}
\toprule
\# & What to do & Effort & Owner? \\
\midrule
\endhead
1 & Make the shipped example run: rename the example leaves to \texttt{A\_a}
    and \texttt{B\_b}, delete \dir{examples/tiny/gene.map}, write the
    underscore rule into \dir{README.md} and
    \dir{examples/tiny/README.md}, and fix the family-listing path bug at
    \loc{gpurec/fit/genewise\_fit.py:79-80} so the command-line route works
    too.  Not closed until the README snippet runs end to end.
  & hours & \\
2 & Fix \texttt{neumann\_terms=3} at
    \loc{gpurec/core/kernels/wave\_backward.py:673}.
  & minutes & \\
3 & Replace the three \texttt{.unwrap()}/\texttt{.expect()} calls in the Rust
    with the crate's own typed error, so a malformed input file raises a normal
    Python exception instead of one that \texttt{except Exception} cannot catch.
  & hours & \\
4 & \textbf{Decide the test story.}  Bring \dir{tests/} onto this branch.  The
    alternative --- deleting the three false claims --- is cheaper but leaves
    every item below unverifiable, so it is not recommended.
  & --- & yes \\
5 & Move \texttt{SolverOptions} out of \dir{gpurec/api/} and
    \texttt{PenaltyOptions} out of \dir{gpurec/solver/}, both into
    \dir{gpurec/config/}.  Removes the deferred-import workaround and both
    layering inversions.  Do this \emph{before} adding any
    \texttt{\_\_init\_\_.py} to \dir{api/} --- see
    Section~\ref{sec:init-trap}.
  & hours & \\
6 & Replace the two \texttt{clamp} calls at
    \loc{gpurec/fit/genewise\_fit.py:328} and \loc{:335} with the
    \texttt{torch.maximum} form already used at
    \loc{gpurec/fit/global\_fit.py:128} and \loc{:136}.
  & minutes & \\
7 & Make \dir{gpurec/fit/newton\_cg.py} read \texttt{NewtonOptions} the way
    \loc{gpurec/fit/genewise\_fit.py:49-56} reads its settings.  Removes 31
    duplicated literals and closes two drifts the code already admits to.
  & hours & \\
8 & Move the six environment variables into settings objects, and give the four
    batch-planning knobs a dataclass alongside \texttt{SolverOptions}.  One of
    them is read at import time to set a GPU launch parameter, so this is not
    purely mechanical.
  & days & \\
9 & Triage the 986 dead lines \emph{with the author}.  456 of them are research
    code that says in a comment it is deliberately not wired in; that is not a
    reviewer's call.
  & hours & yes \\
10 & Restore \dir{docs/latex/kernel\_mathematics.tex}.  Deleting the 16 dangling
    pointers instead is cheaper, but that file is the only thing that ever
    explained the recurrences and eight kernel modules still send readers to it.
  & --- & yes \\
11 & Write a class docstring for \texttt{GeneReconModel} covering its twelve
    constructor arguments, the rate order (duplication, loss, transfer), the
    base-2 logarithm convention, and that the returned loss is in bits.
  & hours & \\
12 & Add READMEs for \dir{solver/}, \dir{fit/}, \dir{config/} and \dir{cli/},
    and add those four directories to the root repository map.
  & hours & \\
13 & Extract the 109-line duplicated Triton block in
    \dir{wave\_backward\_kernels.py} into a device helper, following the
    existing \texttt{\_load\_event\_log\_probability} precedent.
  & hours & \\
14 & Split the three bloated files along the boundaries in
    Section~\ref{sec:bloat}.  Purely structural, nothing depends on it, and
    therefore last.
  & days & \\
\bottomrule
\end{longtable}
\end{center}

\subsection*{The one thing behind most of this}
\label{sec:rootcause}

\texttt{main} is produced by hand-deleting things from \texttt{dev}.  The recent
history reads ``Remove all docs from main'', ``Remove more non-essential docs
from main'', ``Remove planning docs from main'', ``Remove \dir{gpurec/bench}
from main'', ``Remove \dir{gpurec/solver/\_golden} test fixture from main''.

Most of the defects in this review are debris from that process rather than
independent mistakes: the sixteen pointers to deleted documents, the three
false claims that tests exist, the settings guard that cannot fire, the
\dir{notebooks/} directory the README advertises, the pytest configuration with
nothing to configure.  Fixing them one by one leaves the process intact, and the
next strip will produce a fresh set.

\textbf{The durable fix is to generate \texttt{main} from \texttt{dev} with a
script, and to have that script fail when it leaves a reference dangling.}  That
recommendation is not on the numbered list above because it is a change to how
the project is released, not to the code --- but it is worth more than items 10
through 14 combined.

%======================================================================
\section{What \texttt{gpurec} is, and how to read this review}

\subsection{What the software does}

\texttt{gpurec} takes one fixed \emph{species tree} --- the family tree of a set
of organisms --- together with many \emph{gene families}, each being the set of
copies of one gene across those organisms.  For each family it computes how
likely that family's history is, given rates at which genes are duplicated
within a lineage, lost, and transferred sideways from one lineage into another.
Fitting those three rates to the data is what the package is for.  The
arithmetic runs on a graphics card.

The project's own summary, at \loc{README.md:1-8}, calls it ``an experimental
GPU-accelerated likelihood engine for gene-tree/species-tree reconciliation
under a DTL-style model'', which returns ``a scalar negative log-likelihood in
bits''.

Rates can be shared three ways, selected by a \texttt{mode} setting: one set of
rates for everything (\texttt{global}), one per species (\texttt{specieswise}),
or one per gene family (\texttt{genewise}).

\subsection{What this branch is for}

\texttt{main} is the stripped branch.  It carries 96 files; \texttt{dev} carries
those plus 252 more, including 63 test files and 75 documentation files.  The
review assumes \texttt{main} is what an outside user would install, because that
is the only role a branch in this state could serve --- but that assumption is
the reviewer's, not something the repository states, and it is the single
biggest lever on how severe the documentation findings are.  If \texttt{main} is
instead an internal artefact that nobody clones, findings 1 and 2 drop sharply
in importance and little else changes.

\subsection{Words used throughout}

Defining these is not decoration.  One of this review's own findings is that the
codebase uses short names nobody has written down, so the review had better not
do the same thing.

\begin{center}
\captionof{table}{Terms used in this document.  ``Where it comes from'' says
whether the meaning is stated in the repository or was inferred from the code.}
\begin{longtable}{@{}p{0.20\textwidth}p{0.52\textwidth}p{0.20\textwidth}@{}}
\toprule
Term & Meaning & Where it comes from \\
\midrule
\endhead
reconciliation
  & Explaining a gene family's tree as a history inside the species tree: at
    each point, saying whether a gene duplicated, was lost, was transferred, or
    simply followed a speciation.
  & standard usage; the README uses it undefined \\
clade
  & A group of gene copies that could form one subtree.  \texttt{gpurec} works
    over a distribution across possible clades rather than one fixed gene tree.
  & \loc{README.md:15-16} \\
wave
  & One level of the species tree, processed as a batch.  The Rust preprocessor
    groups work into waves so a whole level runs in one kernel launch.
  & \loc{README.md:32-33} \\
the wave step
  & The kernel that advances the per-clade, per-species quantities by one level
    of the species tree.
  & inferred from the code \\
the gene-split reduction
  & Combining the two children of a gene node into their parent, summing over
    the event that joined them.  ``DTS'' in file names.
  & \loc{core/inference/README.md:7} \\
the extinction step
  & Computing, for each species, the probability that a lineage there leaves no
    observed descendants.  Called \texttt{E}; it is a fixed point, solved by
    iterating to convergence.
  & \loc{README.md:35} \\
\texttt{DTS}
  & Duplication, transfer, speciation --- the events combined in a gene-split
    reduction.  Written out in exactly one place in the repository.
  & \loc{core/inference/README.md:7} \\
\texttt{Pi}, \texttt{Pibar}
  & The reconciliation likelihood per clade and species, and its
    transfer-weighted counterpart, propagated wave by wave.
  & \loc{README.md:36}, informally \\
\texttt{E}, \texttt{Ebar}
  & The extinction probabilities above, and their transfer-weighted counterpart.
  & \loc{README.md:35}, informally \\
\texttt{theta}
  & The fitted parameters: base-2 logarithms of the three rates, in the order
    duplication, loss, transfer.
  & \loc{cli/reconcile.py:31} \\
bits vs.\ nats
  & Two units for a log-likelihood: base-2 logarithm versus natural logarithm.
    The library works in bits; the command-line tool prints nats.
  & \loc{cli/\_common.py:8-10} \\
\texttt{vjp}
  & Vector-Jacobian product: the reverse-mode derivative step, i.e.\ what a
    gradient pass computes.  Never expanded in the repository.
  & inferred \\
\texttt{jvp}
  & Jacobian-vector product: the forward-mode derivative.  Never expanded.
  & inferred \\
\texttt{hvp}
  & Hessian-vector product: the second-derivative step.
  & \loc{solver/hvp/\_\_init\_\_.py:1} \\
Newick, ALE
  & Two text formats for trees that the Rust preprocessor reads.
  & \loc{crates/\ldots/lib.rs} \\
\bottomrule
\end{longtable}
\end{center}

\paragraph{The four passes.}  Each of the three pieces of work above exists in
four versions, and the file suffixes encode which: a \emph{value} pass (no
suffix) computes the likelihood; a \emph{gradient} pass
(\texttt{\_backward}) computes its first derivative in reverse mode; a
\emph{forward-derivative} pass (\texttt{\_tangent}) computes a first derivative
in forward mode; and a \emph{second-derivative} pass (\texttt{\_so}, for
``second order'') computes curvature.  All four exist because the package
offers Newton-type fitting, which needs second derivatives, and computing those
efficiently needs the forward-mode pass as an ingredient.  This
three-by-four grid is used later as an argument for how to split two files.

\subsection{Where the three rules come from}
\label{sec:rules}

Three rules are used below to judge the code:

\begin{enumerate}
  \item no default values on function parameters;
  \item no environment variables for configuration;
  \item never clamp a value.
\end{enumerate}

\textbf{These are the project owner's stated working rules, given to the
reviewer directly.  They are not written down anywhere in this repository.}  That
is worth saying plainly, because every other claim in this document carries a
file and a line number and these three cannot.  A reader who does not accept
these rules should read Sections~\ref{sec:stale-default},
\ref{sec:envvars} and~\ref{sec:clamp} as observations rather than as findings
--- although the first of them, the stale default, is a correctness hazard on
any standard.

The reasoning the owner gives for rules 1 and 2 is the same in both cases: one
value should live in one place.  A default is a second home for a number, and it
drifts from the settings file that is supposed to own it, silently, because the
caller that always passes the real value hides the stale default.
Section~\ref{sec:stale-default} is that failure, already happened.

\subsection{What was not reviewed}

This review covers structure, cleanliness, documentation and naming.  It does
\textbf{not} cover: whether the mathematics is correct; whether the numerical
results are accurate; performance; input validation beyond the three Rust
crash sites noted; dependency pinning; or licensing.  No benchmark was run and
no likelihood was checked against a reference implementation.

\subsection{Method}

The branch was checked out into a separate git worktree so the working copy on
\texttt{dev} was untouched.  Four review agents worked in parallel on disjoint
parts of the tree: the GPU kernels and the public model wrapper; the solver,
fitting, command-line and settings code together with the Rust crates; the
documentation; and the naming and layout.

Their findings were then re-checked directly rather than taken on trust.  Every
claim reproduced in Section~\ref{sec:summary}, plus the file and line counts,
the dead-module list, the environment variables, the clamp sites and the import
graph, was verified by re-running the search or the program.  The shipped
example was run through the real Rust preprocessor.  That re-check overturned
two agent claims, both corrected here: the Rust crash sites do not abort the
process (Section~\ref{sec:rust}), and the deferred-import workaround is caused
by \texttt{PenaltyOptions}, not by \texttt{SolverOptions}
(Section~\ref{sec:imports}).  Counts that were not re-derived --- the
duplication measurements and the docstring percentages --- are marked where
they appear.

%======================================================================
\section{How much code is there?}

The branch holds \textbf{24,829 lines across 93 files}, in a tree of 96 tracked
files.  The three not counted are \dir{.gitignore} and the two
\dir{Cargo.lock} files.

\begin{center}
\captionof{table}{Size of the branch.  Line counts are raw lines including
blanks and comments.}
\begin{tabular}{@{}lrr@{}}
\toprule
& Files & Lines \\
\midrule
Python                            & 64 & 18{,}007 \\
Rust                              &  5 &  6{,}149 \\
Markdown (all READMEs)            & 16 &     506 \\
Settings and example data         &  8 &     167 \\
\midrule
\textbf{Counted}                  & \textbf{93} & \textbf{24{,}829} \\
Not counted (\dir{.gitignore}, two \dir{Cargo.lock}) & 3 & --- \\
\midrule
\textbf{Tracked on the branch}    & \textbf{96} & \\
\bottomrule
\end{tabular}
\end{center}

The Python figure of 64 files is 63 under \dir{gpurec/} plus \dir{setup.py} at
the root.  Later sections that scan only the package report 63; both numbers are
used below and are labelled where they appear.

\subsection{There are no tests and no continuous integration}
\label{sec:no-tests}

There is no \dir{tests/} directory, no \texttt{conftest.py}, no \dir{.github/}
directory and no YAML file of any kind on this branch.  Three files say
otherwise.

\begin{center}
\captionof{table}{Claims on this branch that tests exist.}
\begin{tabular}{@{}lp{0.52\textwidth}@{}}
\toprule
Location & Claim \\
\midrule
\loc{pyproject.toml:36}  & \texttt{testpaths = ["tests"]}, plus three test markers \\
\loc{README.md:162}      & tells the reader to run \texttt{pytest} \\
\loc{gpurec/config/defaults.toml:5} & names a ``permanent guard test'' \texttt{tests/test\_config\_toml.py} \\
\bottomrule
\end{tabular}
\end{center}

The third is the damaging one.  That guard is meant to fail if
\dir{defaults.toml} drifts from the Python dataclass it mirrors.  It is not
here, so it cannot fire, and Section~\ref{sec:config-mirrors} shows the drift
it was supposed to catch.

%======================================================================
\section{Are there bloated files?}
\label{sec:bloat}

Three files are bloated in the sense that one file is doing several unrelated
jobs.  Two more are long but cohesive and should be left alone.

\subsection{\texorpdfstring{\dir{crates/gpurec-preprocess/src/lib.rs}}{lib.rs} --- 2{,}401 lines}

Five separate jobs, each already in a contiguous block, so the split is
mechanical.

\begin{center}
\captionof{table}{Proposed split of the largest Rust file.  ``Current lines''
are line ranges in the file as it stands.}
\begin{tabular}{@{}llr p{0.34\textwidth}@{}}
\toprule
Proposed file & Current lines & Size & What it does \\
\midrule
\texttt{ale\_format.rs}   & 922--1480              & 559 & reads the ALE file format \\
\texttt{family\_ccp.rs}   & 1870--2297             & 428 & the \texttt{FamilyCcp} Python class \\
\texttt{amalgamation.rs}  & 757--921, 1678--1869   & 357 & builds clade distributions from gene trees \\
\texttt{newick.rs}        & 153--409               & 257 & parses gene trees in Newick format \\
\texttt{bitset.rs}        & 1580--1677             &  98 & 15 small bit-vector helpers \\
\midrule
\texttt{lib.rs} remainder & ---                    & $\sim$620 & the 4 Python entry points and orchestration \\
\bottomrule
\end{tabular}
\end{center}

There is independent evidence this file was left behind when its siblings were
cleaned up.  The three sibling modules use the crate's typed error type
\textbf{37 times between them and plain string errors zero times}.  This file
still uses plain string errors \textbf{27 times} and the typed error twice.

\subsection{\texorpdfstring{\dir{gpurec/core/kernels/pi\_forward.py}}{pi\_forward.py} --- 1{,}485 lines}

Two jobs: the wave step (lines 91--695 and 1080--1325) and the gene-split
reduction (lines 699--1078 and 1328--1485, about 530 lines).

The codebase itself argues for the split.  In the three-by-four grid described
earlier, the derivative versions of exactly this code are \emph{already}
separated --- \dir{dts\_tangent.py} and \dir{dts\_so.py} hold the gene-split
reduction, \dir{wave\_tangent.py} and \dir{wave\_so.py} hold the wave step.
Only the value pass and the gradient pass keep the two merged.

\subsection{\texorpdfstring{\dir{gpurec/core/kernels/wave\_backward\_kernels.py}}{wave\_backward\_kernels.py} --- 1{,}583 lines}

Three jobs.  The grouping is unambiguous because one can read off which wrapper
in \dir{wave\_backward.py} launches which kernel.

\begin{center}
\captionof{table}{Proposed split of the largest gradient file.  ``Kernel starts
at'' gives the line number where each kernel is defined.}
\begin{tabular}{@{}lrl@{}}
\toprule
Group & Lines & Kernel starts at \\
\midrule
wave-step gradient           & $\sim$883 & 12, 45, 334, 531, 624 \\
gene-split gradient          & $\sim$495 & 906 (one kernel, running to 1400) \\
transfer-complement gradient & $\sim$178 & 1404, 1440, 1480 \\
\bottomrule
\end{tabular}
\end{center}

\subsection{Long but justified --- leave alone}

\dir{crates/gpurec-preprocess/src/scheduler.rs} (1{,}361 lines, 44 functions)
and \dir{batch\_planning.rs} (768 lines, 20 functions) are cohesive: every
function serves one job, wave scheduling and batch packing respectively.

The individual Triton kernels being long is also not a defect.  Triton offers no
cheap function call inside a kernel body, so a 500-line kernel is inherent to
the tool.  The three longest are 495, 458 and 322 lines.

\subsection{The worst single function is not in a bloated file}

\texttt{make\_exact\_hvp\_single} in \loc{gpurec/solver/hvp/exact.py} runs from
line 191 to line 892 --- \textbf{702 lines} --- of which the returned inner
function is \textbf{458 lines in one body} (lines 435--892).  The other four
top-level functions in that 1,000-line file total 148 lines, so the file is
fine and this one function is not.  It already carries three internal section
markers, at lines 372, 817 and 891, at natural split points.  As written it
cannot be tested or profiled piecewise.

Two other Python functions exceed 250 lines: \loc{wave\_backward.py:198} (466
lines) and \loc{api/\_implicit\_grad.py:154} (394 lines).

%======================================================================
\section{Is the code clean?}

Inside the source files, yes.  The problems are in how settings and dead code
are managed, not in how the code is written.

\subsection{What came back clean}

Zero \texttt{TODO}, \texttt{FIXME}, \texttt{XXX} or \texttt{HACK} markers
anywhere; zero bare \texttt{except:} clauses; zero mutable default arguments;
zero commented-out code; zero stray debug prints (all 55 \texttt{print()} calls
are behind a verbosity flag, are a command-line message, or go to standard
error); and in the Rust, zero \texttt{unsafe} blocks and zero explicit
\texttt{panic!}, \texttt{unreachable!}, \texttt{todo!} or
\texttt{unimplemented!} macros across 6{,}149 lines.  Nesting nowhere exceeds
five levels.

That is a better result than most codebases return, and it is why the verdict
separates the code from its surroundings.

\subsection{Three Rust call sites raise an exception nobody can catch}
\label{sec:rust}

Three places in the Rust use \texttt{.unwrap()} or \texttt{.expect()} on a value
that a malformed input file can make absent:

\begin{lstlisting}
crates/gpurec-preprocess/src/lib.rs:1109
    let observations = ale.observations.expect("validated observations");
crates/gpurec-preprocess/src/lib.rs:1806
    let idx = *split_index_map.get(&canon).unwrap();
crates/gpurec-backtrack/src/lib.rs:413
    let sampler = self.recipient_cache[clade].as_ref().unwrap();
\end{lstlisting}

All three are reachable from Python.  \textbf{What happens is not a crash of the
process}, and an earlier draft of this review said it was.  Neither
\dir{Cargo.toml} sets \texttt{panic = "abort"}, so the default unwinding
behaviour applies and the Python binding layer catches the panic at the
boundary.  The user sees \texttt{pyo3\_runtime.PanicException}.

That is still a real defect, for two reasons.  First, that exception derives
from Python's \texttt{BaseException} rather than \texttt{Exception}, so an
ordinary \texttt{try/except Exception} around a file-loading call will not catch
it --- a caller processing a thousand families cannot skip the one bad file and
continue.  Second, it bypasses the crate's own well-built typed error, which
already knows how to turn a failure into the right Python exception class.  The
fix is to use that error type, and it is item 3 on the list.

\subsection{986 lines have no caller}
\label{sec:dead-code}

Each was checked by searching the whole tree --- Python, Rust, Markdown and
plain text --- for any reference outside the defining file.

\begin{center}
\captionof{table}{Code with no caller anywhere in the repository.}
\begin{tabular}{@{}lr@{}}
\toprule
Item & Lines \\
\midrule
\dir{gpurec/solver/curvature/genewise.py} --- the whole file & 456 \\
\dir{gpurec/distributed.py} --- the whole file               & 112 \\
\dir{gpurec/fit/baselines.py} --- the whole file             &  62 \\
\texttt{newton\_tr} and \texttt{newton\_cg}, \loc{fit/newton\_cg.py:274-435} & 162 \\
\texttt{receiver\_information}, \loc{solver/curvature/receiver.py:195}       &  75 \\
\texttt{origination\_information}, \loc{solver/curvature/origination.py:143} &  70 \\
\texttt{plan\_family\_batches}, \loc{core/scheduling/batching.py:394}        &  49 \\
\midrule
\textbf{Total} & \textbf{986} \\
\bottomrule
\end{tabular}
\end{center}

The first entry says so itself: \loc{solver/curvature/genewise.py:60} carries
the comment ``TEST-ONLY / NOT wired into the fit \ldots do NOT wire it into the
fit''.  That is 456 of the 986 lines, and it is why item 9 on the list says to
triage this with the author rather than delete it.

A distinction worth keeping: a further body of code is unreachable from the two
entry points but is \emph{intentional} research surface.
\loc{gpurec/fit/dtl\_fit.py:73} deliberately raises \texttt{NotImplementedError}
for one fitting mode and names the hand-import entry points in the error
message.  That is the right pattern.  What the 986 lines above lack is either a
caller or such a pointer.

\subsection{A stale default that is 21 times too small}
\label{sec:stale-default}

\loc{gpurec/core/kernels/wave\_backward.py:673} declares
\texttt{neumann\_terms=3}.  The configured value is \textbf{64}, in two places:
\loc{gpurec/api/solver\_options.py:12} and
\loc{gpurec/config/defaults.toml:22}.

That setting controls how many terms of a series are summed when solving for
the gradient --- step 7 of the README's execution path, ``a matrix-free
Neumann-series solve for the E-adjoint linear system''.  Both production callers
pass the value explicitly, which is exactly why the stale \texttt{3} has gone
unnoticed: it never runs in production.  It runs for whoever calls that function
by hand and gets a gradient built from 3 terms instead of 64, with no warning.

\textbf{How wrong is the answer?}  This review does not know, and that is a real
gap.  The size of the error depends on how fast that series converges for this
model, which was not measured --- measuring it was outside the scope stated
above.  What can be said without measuring: the project chose 64 and wrote it
into the settings file twice, so 3 is not a value anyone deliberately selected,
and a silent 21-fold reduction in solver work is not something a caller should
receive by accident.  Anyone fixing item 2 should take ten minutes to compare
the two gradients on a small family and record the answer, because it converts
this from a hazard into a known quantity.

A second instance of the same shape: \loc{gpurec/api/\_implicit\_grad.py:82}
hardcodes \texttt{max\_iter: int = 128}, restating a value already declared at
\loc{gpurec/api/solver\_options.py:17}.

\paragraph{The measurement.}  Across \dir{solver/}, \dir{fit/}, \dir{cli/},
\dir{optimization.py} and \dir{distributed.py} there are \textbf{446 defaults on
817 parameters, or 55\,\%}.  Of those 446: \textbf{219} are \texttt{None}
sentinels, \textbf{142} are bare numeric literals, 44 are booleans, 20 are
strings and 21 are other expressions.

The 142 numeric literals are the ones that matter, because a number is what
drifts.  The \texttt{None} sentinels are a different thing wearing the same
clothes: most of them mean ``the caller did not supply this'', and several are
the good pattern where the function then reads the real value from the settings
object --- one value, one place, which is what the rule is protecting.  The
booleans, strings and other expressions were not individually assessed.

The worst offenders by numeric literal are \dir{fit/newton\_cg.py} (31),
\dir{fit/optimize.py} (28), \dir{fit/global\_fit.py} (15) and
\dir{fit/genewise\_fit.py} (13).

The rule \emph{is} followed in one place and it works:
\loc{gpurec/fit/specieswise\_fit.py:29} makes the regularisation weight a
required argument and raises a clear error explaining that a weight of zero is
the unregularised fit and is deliberately not a default.

\subsection{The same settings in three or four places}
\label{sec:config-mirrors}

\dir{gpurec/config/defaults.toml} is \textbf{never read at runtime} unless the
user passes \texttt{-{}-config}.  \loc{gpurec/config/gpurec\_config.py:210}
returns the plain Python dataclass values with no file access.  So the file is a
hand-maintained copy of 45 numbers with nothing checking it --- the guard test
it names is not on this branch.

The Newton solver's settings exist three times over: as 21 fields in
\dir{gpurec/config/newton.py}; repeated at \loc{defaults.toml:39-59}; and
re-typed as function-parameter defaults at \loc{fit/newton\_cg.py:49-52},
\loc{:274-276} and \loc{:366-368} --- a file containing \textbf{zero} references
to the settings class.

The code already admits the drift.  \loc{gpurec/config/newton.py:35-39} says two
fields are ``RESERVED: validated here but not yet read by
\dir{fit/newton\_cg.py} (FD literal is inline)''.
\loc{gpurec/fit/optimize.py:517-519} notes the same setting standing at 8 in one
place and 40 in another.

A fourth partial mirror sits at \loc{gpurec/config/gpurec\_config.py:159}
and~\loc{:175}, re-typing solver values already identical to the dataclass
defaults.  And four batch-planning settings have \textbf{no settings home at
all}, appearing as defaults in five signatures:
\texttt{family\_chunk\_size=300}, \texttt{clade\_budget=315{,}000},
\texttt{batch\_packing="depth\_first\_fit"}, \texttt{max\_wave\_size=8192}.

Credit where due: \loc{fit/genewise\_fit.py:49-56} and \loc{fit/map\_cv.py:46}
\emph{derive} their solver settings from the factory methods instead of
retyping them.  That is the pattern item 7 should copy.

\subsection{Six files take settings from the environment}
\label{sec:envvars}

\begin{center}
\captionof{table}{Environment variables used as settings.}
\begin{longtable}{@{}p{0.40\textwidth}p{0.54\textwidth}@{}}
\toprule
Location & What it does \\
\midrule
\endhead
\loc{gpurec/core/memory\_policy.py:69-70}
  & \textbf{The worst two.}  Read as
    \texttt{os.environ.get(NAME, str(default\_fraction))} --- the shell variable
    is the first choice and the argument the caller explicitly passed is only
    the fallback.  A deliberate setting is silently overridden. \\
\loc{gpurec/fit/genewise\_fit.py:245-374}
  & \emph{Writes} a variable before the fit, pops it mid-loop to force a cold
    re-check, restores it in a \texttt{finally}.  Nine statements. \\
\loc{gpurec/core/kernels/wave\_tangent.py:24}
  & A GPU launch parameter read at import time, so it cannot appear in a
    settings file or in the record of a run. \\
\loc{gpurec/solver/hvp/exact.py:207, 233}
  & Two numeric solver knobs. \\
\loc{gpurec/api/\_execution.py:138, 310} and \loc{solver/hvp/exact.py:62}
  & Switch a real behaviour: whether the previous call's gradient is reused as
    a starting guess. \\
\loc{gpurec/distributed.py:37-41}
  & Multi-process settings.  (This file is dead; see
    Section~\ref{sec:dead-code}.) \\
\bottomrule
\end{longtable}
\end{center}

Two observations.  \loc{gpurec/api/solver\_options.py:27-34} documents a sibling
setting and states ``config-only, no env var'' --- so this is a known migration
that stalled partway.  And \loc{gpurec/config/memory.py:28-29} documents the
memory variables in comments while the dataclass field beside them holds the
same value, which is two homes for one setting.

One use is \emph{not} a violation and should be kept:
\loc{gpurec/core/native.py:15} reads a path variable to locate the compiled Rust
library, and documents it in the error message at line 53.  That is a fact about
where a build artefact lives, not a setting.

The Rust uses no environment variables.

\subsection{Clamping}
\label{sec:clamp}

Two real violations, and in both cases the compliant form already exists two
files away:

\begin{itemize}
  \item \loc{gpurec/fit/genewise\_fit.py:328} --- \texttt{e.clamp(min=mu)} on
        curvature eigenvalues.  \loc{gpurec/fit/global\_fit.py:128} does the
        same job with \texttt{torch.maximum}.
  \item \loc{gpurec/fit/genewise\_fit.py:335} ---
        \texttt{dn.clamp(min=trust)}.  \loc{gpurec/fit/global\_fit.py:136} uses
        \texttt{torch.maximum}.
\end{itemize}

One case needs a decision rather than a fix.  \loc{gpurec/optimization.py:95}
(\texttt{clamp\_log\_rate\_}) clamps the fitted rates onto their allowed range.
It is deliberate, it is public API, and five call sites depend on it.  The
options are: keep it and record in a comment that this is the one sanctioned
clamp; or replace it with a projection that reports when it binds, so a rate
pinned at its bound is visible rather than silent.  The second is
better if a bound-hitting rate would otherwise be mistaken for a converged fit.

Five further uses were checked and are legitimate: an array-index bound at
\loc{core/parameters/extract\_parameters.py:137} (whose comment records that the
old value-clamp was deliberately removed); the denominator of a reported
diagnostic at \loc{wave\_backward.py:535}; and three guards on a variance before
a square root, in \dir{solver/curvature/}.

\subsection{Duplication}

Measured by the review agents; the two largest were re-checked directly.

\paragraph{109 byte-identical lines in one file.}  In
\dir{gpurec/core/kernels/wave\_backward\_kernels.py}, the block at lines
156--264 and the block at lines 723--831 are byte-for-byte identical.  Counting
shorter runs, the two kernels at lines 45 and 624 share \textbf{156 identical
lines} across four runs; both kernels are about 280 lines, so this is over half
of each.  The duplicated block is shared setup: loading the event constants,
gathering the two child rows, forming the loss and transfer terms.

This is clearly fixable, because the project already does it elsewhere:
\dir{pi\_forward.py} defines shared device helpers and two other kernel files
import one across file boundaries.

\paragraph{About a third of a 476-line file is a copy of another part of it.}  In
\dir{gpurec/api/\_execution.py}, two functions of 102 and 104 lines share
\textbf{94 identical lines}, differing in exactly four things: the name, an extra
flag, an up-front error check, and which of two loss functions they call.  One
level up, two more functions share 42 of about 55 lines.  Together that is
roughly 150 of the file's 476 lines.  The environment-variable block noted
earlier exists in duplicate at lines 138 and 310 for this reason.

\paragraph{178 near-copied lines across three curvature files.}  Counting a pair
of lines as a copy when they are at least 70\,\% character-identical after
normalising whitespace --- a fuzzy measure, so treat these as indicative:

\begin{center}
\captionof{table}{Near-copies within \dir{solver/curvature/}.  ``Share'' is of
the smaller of the two function bodies.}
\small
\begin{tabular}{@{}p{0.56\textwidth}rr@{}}
\toprule
Pair & Lines & Share \\
\midrule
\texttt{receiver.py:273} vs \texttt{origination.py:216} & 50 & 85\,\% \\
\texttt{receiver.py:195} vs \texttt{origination.py:143} & 45 & 73\,\% \\
\texttt{origination.py:216} vs \texttt{genewise.py:223} & 36 & 60\,\% \\
\texttt{receiver.py:147} vs \texttt{origination.py:108} & 26 & 87\,\% \\
\texttt{origination.py:108} vs \texttt{genewise.py:184} & 21 & 70\,\% \\
\midrule
\textbf{Total} & \textbf{178} & \\
\bottomrule
\end{tabular}
\end{center}

Those three files total 1{,}083 lines, so this is about 16\,\%.  Because
\dir{origination.py} appears in three rows, the 178 counts each \emph{pairing},
not each distinct line --- the number of lines that could be deleted is smaller
than 178 and was not measured.  A sixth pairing, of 9 lines, sits between
\dir{global\_fit.py} and \dir{genewise\_fit.py} and is discussed next.

Credit where due: \dir{solver/curvature/gauge.py} already holds the shared core,
and \loc{origination.py:38} imports helpers from \dir{receiver.py} rather than
copying them.  One round of de-duplication has happened; the wrapper layer is
what remains.

\paragraph{One of these is a behaviour difference, not repetition.}
\loc{fit/genewise\_fit.py:326} builds the small curvature matrix analytically,
while \loc{fit/global\_fit.py:120-125} rebuilds the same matrix by finite
differences --- three extra forward-and-backward passes per refresh --- even
though its own docstring says it runs ``the SAME recipe''.  Worth resolving
independently of the duplication.

\paragraph{Two more pairs in the kernels.}
\loc{wave\_tangent.py:563-648} and \loc{:650-734} are two 85-line launchers with
73 lines identical.  \loc{:28-272} and \loc{:276-487} share 130 identical lines,
the second being the first wrapped in a repeat loop.

\paragraph{What is \emph{not} duplicated.}  The four passes are genuinely
different code, not variants of one thing.  The three \texttt{e\_step} files
share only their input loading; the mathematics differs.  The real duplication
is \emph{within} files, not across the passes.

\subsection{Smaller items}

\begin{itemize}
  \item \textbf{Nine dead defensive lookups} of the form
        \texttt{getattr(static, "field", fallback)} --- at
        \loc{api/\_execution.py:23, 73, 138, 310}, \loc{api/\_autograd.py:36, 70},
        \loc{core/inference/solver.py:42, 85} and
        \loc{core/backtracking/input.py:57} --- all name fields that
        \emph{are} declared on the dataclass at \loc{api/\_batch\_state.py:13-36}.
        The fallback can never fire, and it is worse than nothing: a mistyped
        field name would silently return the fallback instead of raising.
  \item \textbf{A hardcoded personal path in library code.}
        \loc{gpurec/fit/map\_cv.py:173} contains an absolute path under one
        developer's home directory, inside a self-test wired to a
        \texttt{\_\_main\_\_} block.  That directory does not exist even in this
        checkout.
  \item \textbf{One raw \texttt{print} to standard output} at
        \loc{gpurec/api/model.py:177-182}.  The information is legitimate; it
        belongs in a warning or a logger so callers can silence or capture it.
  \item \textbf{A magic iteration count with no home}, \texttt{400} at
        \loc{gpurec/api/model.py:366}, with no comment explaining it.
  \item Two stale path comments: \loc{gpurec/solver/penalties.py:1} names a
        directory that does not exist, and \loc{gpurec/config/newton.py:4}
        points at a file that became a package in commit \texttt{5a5b0f33}.
\end{itemize}

%======================================================================
\section{Is the documentation good?}

No.  Two problems dominate: the example does not run, and sixteen pointers lead
to deleted files.

\subsection{The shipped example fails, and the rule it breaks is written nowhere}
\label{sec:docs-example}

Run against the real Rust preprocessor:

\begin{lstlisting}
preprocess_dataset(examples/tiny/species.nwk, [examples/tiny/gene.nwk])
  -> ValueError: species "a" not found for gene leaf "a"
\end{lstlisting}

\texttt{gpurec} maps a gene to its species by splitting the leaf name on the
\textbf{first underscore} --- \texttt{split\_once('\_')} at
\loc{crates/gpurec-preprocess/src/lib.rs:516}, and again at line 1657.  That is
the only mechanism and it is documented nowhere: not in a README, not in a
docstring, not in a comment.

\dir{examples/tiny/gene.nwk} contains \texttt{(a:1,b:1);} with no underscores,
so the headline example at \loc{README.md:47-66} fails on the only dataset the
repository ships.  Renaming the leaves to \texttt{A\_a} and \texttt{B\_b} was
tested and gets past the leaf-mapping stage; \textbf{whether the full README
snippet then runs to completion was not confirmed}, which is why item 1 on the
list is not closed until someone runs it end to end.

Two further problems in the same example:

\begin{itemize}
  \item \dir{examples/tiny/gene.map}, containing \texttt{A:a} and \texttt{B:b},
        is a decoy.  Searching every Python and Rust file finds \textbf{nothing
        that reads it}, yet \loc{examples/README.md:13-14} and
        \loc{examples/tiny/README.md:12-13} both present it as the thing that
        assigns genes to species.  A new user will edit that file and wonder why
        nothing changes.
  \item \dir{examples/tiny/families.txt} names its gene tree as
        \texttt{gene.nwk}, and the resolver at
        \loc{gpurec/fit/genewise\_fit.py:79-80} takes that string verbatim
        without joining it to the listing file's own directory.  So it resolves
        relative to the current working directory, and the command-line route
        into the same example cannot find the file from the repository root.
\end{itemize}

\subsection{Sixteen pointers to deleted documentation}

All of \dir{docs/} was removed from this branch; the references to it were not.

\begin{center}
\captionof{table}{References to documentation files that are not on this branch.}
\begin{tabular}{@{}p{0.30\textwidth}rp{0.50\textwidth}@{}}
\toprule
Missing target & Refs & Referenced from \\
\midrule
\dir{docs/latex/kernel\_mathematics.tex} & 9
  & \loc{core/kernels/README.md:8} (a clickable link), plus line~1 of six kernel
    modules, \loc{dts\_tangent.py:191} and \loc{wave\_backward.py:3} \\
\dir{docs/config\_convention.md} & 3
  & \loc{fit/optimize.py:484}, \loc{fit/map\_cv.py:77}, \loc{fit/genewise\_fit.py:60} \\
\dir{docs/pi\_storage.md} & 1 & \loc{api/README.md:12} \\
\dir{docs/centered\_state\_contract.md} & 1 & \loc{core/pi\_state.py:4} \\
\dir{docs/superpowers/\ldots} & 2
  & \loc{core/parameters/fraction\_missing.py:6}, \loc{fit/specieswise\_fit.py:5} \\
\bottomrule
\end{tabular}
\end{center}

The first row is why item 10 recommends restoring rather than deleting.
\loc{gpurec/core/kernels/README.md:8-9} states: \emph{``The canonical equations
\ldots are in [that file].  Python docstrings are intentionally limited to
interface and launch behavior.''}  That was a sound policy while the file
existed.  Deleting it turned the policy into a hole --- six kernel modules now
carry one-line docstrings whose whole content is ``see this file''.

Three more dead pointers sit in \dir{pyproject.toml}: a plan file at line 12, an
experiments script at line 25, and \dir{solver/cg.py} at line 17 (the file is
\dir{solver/krylov.py}).

\subsection{Docstring coverage}

Measured by a review agent over all 64 Python files.  ``Public functions'' means
functions and methods whose name does not begin with an underscore.

\begin{center}
\captionof{table}{Docstring coverage across the branch.}
\begin{tabular}{@{}lrr@{}}
\toprule
& Documented & Share \\
\midrule
Module docstrings            & 43 of 64   & 67\,\% \\
Public functions and methods & 126 of 239 & 53\,\% \\
Private functions            & 58 of 119  & 49\,\% \\
Classes                      & 11 of 17   & 65\,\% \\
\bottomrule
\end{tabular}
\end{center}

Coverage is sharply split and \emph{inversely} correlated with README coverage:
the directories with no README are the best documented in code, and the
directories with READMEs are the worst.

\begin{center}
\captionof{table}{Best and worst documented files.  Fractions count public
functions only, so a large kernel file can show a small denominator: the Triton
kernels themselves are private, and only the Python launch wrappers are counted.}
\begin{tabular}{@{}p{0.42\textwidth}p{0.50\textwidth}@{}}
\toprule
Well documented & Barely documented \\
\midrule
\loc{config/gpurec\_config.py} 10 of 10 &
  \loc{core/kernels/pi\_forward.py} \textbf{0 of 3}, in 1{,}485 lines \\
\loc{solver/curvature/receiver.py} 8 of 10 &
  \loc{core/kernels/e\_step.py} \textbf{0 of 5}, in 605 lines \\
\loc{solver/curvature/origination.py} 7 of 8 &
  \loc{core/inference/forward.py} \textbf{0 of 1} --- the forward pass itself \\
\loc{solver/penalties.py} 8 of 10 &
  \loc{api/\_implicit\_grad.py} \textbf{0 of 2}, in 670 lines \\
\loc{fit/optimize.py} 5 of 6 &
  \loc{api/\_autograd.py} \textbf{0 of 4} functions, 0 of 2 classes \\
\bottomrule
\end{tabular}
\end{center}

\paragraph{The main public class has no docstring at all.}
\texttt{GeneReconModel} at \loc{gpurec/api/model.py:34} --- the one object the
README tells the reader to use --- has no class docstring, no constructor
docstring, no \texttt{forward} docstring, and its module has none either.  Its
constructor takes twelve keyword arguments, three of which are unexplained
magic numbers.  The three methods the README advertises are all undocumented.
That is item 11.

\paragraph{A concrete bad example.}  \loc{core/kernels/pi\_forward.py:1156} is a
public launch wrapper with \textbf{24 parameters}, no docstring, no type hints
and no shape annotations.  A reader cannot tell what any tensor's shape is
without tracing a caller.  Two neighbouring wrappers have 21 and 22 parameters.

\paragraph{Two good examples worth copying.}  \loc{fit/global\_fit.py:1-19}
states the parameter vector and its shape, the objective, the recipe, \emph{why}
families are never dropped in that mode, and the measured speedup.
\loc{solver/curvature/receiver.py:1-30} states the invariance, why it makes the
curvature matrix singular, and what projection removes it.

\subsection{Comments in the kernels explain the wrong layer}

Comment density in \dir{core/kernels/} runs 0--6\,\%, against 13\,\% in
\dir{solver/hvp/exact.py}.  The bare numbers are less interesting than the
pattern, so the table is omitted; what matters is \emph{what} is commented.

The comments that exist are good, and they cover numerical and hardware
concerns: why an index is widened to 64-bit to avoid overflow
(\loc{pi\_forward.py:227-228} and two others), why a row offset is absorbed
lazily (\loc{pi\_forward.py:441-443}), and an excellent five-line explanation of
a cross-thread lost-update race at \loc{e\_step.py:216-220}.

What is absent is the model.  Nowhere in the kernel directory is there a
statement of what the extinction probability or the reconciliation quantities
are, or what the recurrence computes.  The fixed point at the heart of the
model, \loc{e\_step.py:19}, has no docstring and two comments, neither about the
recurrence.

Two mitigations: the names are unusually self-documenting, and
\loc{wave\_backward\_kernels.py:912-945} annotates each of about 30 pointer
arguments inline with its shape and meaning.  That block is the best tensor
documentation on the branch and is the right template for the
\dir{pi\_forward.py} wrappers.

Note that density alone is not the standard here.  \dir{solver/hvp/exact.py} has
the highest density on the branch and also contains the 702-line function
criticised in Section~\ref{sec:bloat}; good comments do not excuse structure,
and the point is that the kernels lack the \emph{kind} of comment that would
help, not that they lack volume.

\subsection{The READMEs}

Sixteen files, 506 lines.  Five are genuine and useful: the root one,
\dir{crates/README.md}, \dir{crates/gpurec-preprocess/src/README.md},
\dir{gpurec/core/README.md}, and \dir{gpurec/core/kernels/README.md} --- the
best of them.  Four are nine-line stubs that restate the directory name and end
with the same filler sentence about \texttt{\_\_pycache\_\_} being an
interpreter artefact.

\begin{center}
\captionof{table}{README coverage against size.  Line counts are Python only.}
\begin{tabular}{@{}lrrl@{}}
\toprule
Directory & Files & Lines & READMEs \\
\midrule
\dir{gpurec/core/}   & 21 & 9{,}554 & 6 \\
\dir{gpurec/solver/} & 13 & 3{,}572 & \textbf{none} \\
\dir{gpurec/fit/}    &  9 & 1{,}942 & \textbf{none} \\
\dir{gpurec/api/}    &  6 & 1{,}942 & 1 \\
\dir{gpurec/config/} &  6 &   470 & \textbf{none} \\
\dir{gpurec/cli/}    &  5 &   291 & \textbf{none} \\
\bottomrule
\end{tabular}
\end{center}

The root README's repository map at \loc{README.md:121-134} never mentions
\dir{fit/}, \dir{solver/}, \dir{cli/} or \dir{config/}.

\paragraph{Stale claims in the root README.}  The headline example fails
(lines 47--66); a solver setting \texttt{e\_init} is documented (line 111) that
does not exist; a \dir{notebooks/} directory is advertised (line 134) that is
not here; the dependency list (line 138) omits scipy, which
\loc{pyproject.toml:20} requires; external reference projects are described
(lines 158--160) that do not exist; and the \texttt{pytest} advice
(lines 162--164) concerns tests and vendored directories that are both absent.
Separately \loc{gpurec/README.md:9-20} lists six public exports where there are
seven.

\subsection{The command-line tool is undocumented}

\loc{pyproject.toml:29} installs a \texttt{gpurec} command with two
subcommands, \texttt{reconcile} and \texttt{fit}.  No README mentions it.  Its
input formats --- a directory, a glob, or a family listing file --- and its
output format are findable only by reading source, as is the fact that it prints
in nats while the library works in bits.

\subsection{What a new user cannot work out}

In rough order of how badly it blocks them: the gene leaf naming rule; that
\dir{gene.map} does nothing; the order of the three rates (duplication, loss,
transfer --- stated exactly twice in the whole repository, at
\loc{cli/reconcile.py:31} and \loc{fit/global\_fit.py:3}, so anyone reading the
parameters directly will guess wrong five times in six); that the parameters are
base-2 logarithms; that the loss is in bits; the entire command-line surface;
that a \texttt{-{}-config} file exists at all; the three undocumented
environment variables; and the vocabulary, for which this document's own
glossary is offered as a starting point.

%======================================================================
\section{Are the naming conventions consistent?}

\textbf{Yes, mechanically --- this is not a problem area.}  A syntax-tree scan
over the 63 Python files under \dir{gpurec/} found zero classes not in CapWords,
zero module-level constants not in upper case, and only 15 of 358 functions not
in lower-case-with-underscores --- 13 of those being deliberate mathematical
shorthand for small local closures.  Cross-file naming for the same concept is
consistent: the species tree is called \texttt{species\_tree} in all 18 places
and never anything else.

There is one real problem, and it is the same one the review's own glossary
exists to answer.

\subsection{Short names nobody has written down}

\begin{center}
\captionof{table}{Frequently-used short names and where, if anywhere, the
repository defines them.}
\begin{tabular}{@{}p{0.16\textwidth}p{0.28\textwidth}rp{0.34\textwidth}@{}}
\toprule
Name & Meaning & Uses & Defined in the repository \\
\midrule
\texttt{S}  & number of species        & 114 & once, \loc{wave\_backward.py:706} \\
\texttt{W}  & wave width               &  22 & \textbf{nowhere} \\
\texttt{G}  & number of families       &  17 & \textbf{nowhere} \\
\texttt{C}  & number of clades         &  11 & \textbf{nowhere} \\
\texttt{sv} & forward-solve results    &  41 & \textbf{nowhere} (called \texttt{saved} in one docstring) \\
\texttt{Pi}, \texttt{Pibar}, \texttt{E}, \texttt{Ebar} & see glossary & 700+ & informally, in the README's prose \\
\texttt{DTS} & duplication, transfer, speciation & 35 & once, \loc{core/inference/README.md:7} \\
\texttt{vjp}, \texttt{jvp} & see glossary & many & \textbf{nowhere} \\
\texttt{th}, \texttt{al}, \texttt{gp}, \texttt{c1}, \texttt{c2}, \texttt{lam} & --- & 110 & \textbf{nowhere} \\
\bottomrule
\end{tabular}
\end{center}

On top of these are invented step labels with no key: \texttt{S3} through
\texttt{S9} appear in ten comments across \dir{solver/hvp/}, and
\loc{solver/curvature/receiver.py:1} opens with ``S9: gauge-projected
joint\ldots''.  None is defined anywhere.  Likewise
\texttt{proposal0\_memory\_gate}, exported from \dir{core/memory\_policy.py} ---
``proposal0'' appears nowhere else.

Three abbreviations are handled correctly and show the authors know how:
``second order'' is spelled out in line 1 of each \texttt{\_so} file,
``Hessian-vector product'' in \loc{solver/hvp/\_\_init\_\_.py:1}, and
``Gauss-Newton'' in \loc{solver/hvp/gauss\_newton.py:1}.

One collision: \texttt{so} means ``second order'' in file names but is a local
variable holding a settings object in three files.

\subsection{Two names that mislead}

\dir{gpurec/optimization.py} versus \dir{gpurec/fit/optimize.py}: the first is
rate-bound projection (96 lines), the second is the optimiser loop (561 lines),
and nothing in either name distinguishes them.

The \texttt{*\_fit.py} suffix covers four of the eight fitting modules but not
\dir{baselines.py}, \dir{map\_cv.py}, \dir{newton\_cg.py} or
\dir{optimize.py} --- and \dir{dtl\_fit.py} breaks the pattern's reading, being
the dispatcher across modes rather than a mode.

%======================================================================
\section{Is the directory tree well organised?}
\label{sec:imports}

The shape is good.  Two settings classes are in the wrong package, and that
alone explains every layering problem below.

\subsection{The intended layering, and where it is violated}

Reading the package descriptions, the intended stack is:

\begin{center}
\texttt{config} $\;\rightarrow\;$ \texttt{core} $\;\rightarrow\;$
\texttt{solver} $\;\rightarrow\;$ \texttt{api} $\;\rightarrow\;$
\texttt{fit} $\;\rightarrow\;$ \texttt{cli}
\end{center}

reading left to right as ``may be imported by''.  Settings sit at the bottom and
know about nothing; the command-line tool sits at the top and may use anything.
An arrow pointing right-to-left is a violation.

\begin{center}
\captionof{table}{Imports between packages.  Numbers are counts of import
statements; imports within a package are excluded.  A dagger marks a direction
that violates the layering above.}
\begin{tabular}{@{}ll@{}}
\toprule
Package & Imports from \\
\midrule
\texttt{api}    & \texttt{core} 15, \texttt{config} 4 \\
\texttt{cli}    & \texttt{config} 4, \texttt{fit} 2, \texttt{api} 1 \\
\texttt{config} & \texttt{solver} 4$^\dagger$, \texttt{api} 2$^\dagger$ \\
\texttt{core}   & \texttt{config} 2, \texttt{api} 2$^\dagger$ \\
\texttt{fit}    & \texttt{solver} 17, \texttt{config} 7, \texttt{api} 5, \texttt{core} 1 \\
\texttt{solver} & \texttt{core} 26, \texttt{config} 8, \texttt{api} 5$^\dagger$ \\
\bottomrule
\end{tabular}
\end{center}

Every violating edge is one of two imports:

\begin{itemize}
  \item \texttt{config} $\rightarrow$ \texttt{api} (2 statements) and
        \texttt{core} $\rightarrow$ \texttt{api} (2 statements) are
        \emph{all four} the same import: \texttt{SolverOptions} from
        \dir{gpurec/api/solver\_options.py}.  That file is 58 lines --- one
        dataclass and a validation method --- with \textbf{zero} imports from
        \texttt{gpurec}.  It is pure settings sitting in the model package.
  \item \texttt{config} $\rightarrow$ \texttt{solver} (4 statements, two of
        them comments) is entirely \texttt{PenaltyOptions} from
        \dir{gpurec/solver/penalties.py} --- again a settings dataclass, this
        time in the solver package.
\end{itemize}

Moving both into \dir{gpurec/config/} removes every violating edge on the
branch.  That is item 5.

\subsection{The workaround, and a trap in removing it}
\label{sec:init-trap}

\texttt{PenaltyOptions} causes a genuine circular import.  The settings module
cannot import it at the top of the file, because doing so pulls in the solver
package, which pulls in the model package, which imports the settings module
again while it is still half-built.  The workaround is to defer the import
inside a function, and it carries a nineteen-line comment at
\loc{gpurec/config/gpurec\_config.py:21-39} explaining exactly this.

\textbf{An earlier draft of this review attributed that comment to
\texttt{SolverOptions}.  That was wrong, and the comment itself says so}, ending:
``\texttt{SolverOptions} has no such cycle (\texttt{gpurec.api} is a plain
namespace package with no \texttt{\_\_init\_\_.py}, and \texttt{solver\_options.py}
has no gpurec imports of its own), so it stays a normal top-level import.''

That sentence is worth reading twice, because it sets a trap.
Section~\ref{sec:init} notes that six package directories have no
\texttt{\_\_init\_\_.py} and calls it an inconsistency.  But the missing
\texttt{\_\_init\_\_.py} in \dir{api/} is precisely what stops the
\texttt{SolverOptions} edge from becoming a real cycle.  \textbf{Adding an
\texttt{\_\_init\_\_.py} to \dir{api/} --- exactly the tidying that
Section~\ref{sec:init} invites --- would create a second circular import},
unless \texttt{SolverOptions} is moved out first.  Hence the ordering note on
item 5.

\subsection{Six package directories have no \texttt{\_\_init\_\_.py}}
\label{sec:init}

They are \dir{gpurec/api/}, \dir{gpurec/core/}, and the four \dir{core/}
subpackages.

This is \textbf{not} a packaging break.  A wheel was built to check: all six
subpackages ship (68 entries) and both
\texttt{import gpurec.api.model} and
\texttt{import gpurec.core.kernels.pi\_forward} succeed, because the packaging
tool discovers namespace packages by default.

The cost is intent.  \dir{solver/} and \dir{config/} each declare an export
list; \dir{api/}, which the root README calls ``the stable interface'', declares
nothing, so there is no single place saying what it offers.  Given
Section~\ref{sec:init-trap}, the right response is not to add the files but to
move the settings classes first and then decide.

\subsection{The private-module convention is not enforced}

\dir{gpurec/api/} marks four modules private with a leading underscore.  Six
files outside \dir{api/} import from them:

\begin{lstlisting}
gpurec/solver/value_and_grad.py:24    from gpurec.api._batch_state   import _BatchStatic
gpurec/solver/value_and_grad.py:25    from gpurec.api._execution     import stream_batches
gpurec/solver/curvature/genewise.py:6 from gpurec.api._execution     import evaluate_static_loss_grad, stream_batches
gpurec/solver/hvp/exact.py:18         from gpurec.api._implicit_grad import _neumann_e_adjoint, _safe_exp2_ratio
gpurec/solver/hvp/gauss_newton.py:22  from gpurec.api._implicit_grad import _safe_exp2_ratio, implicit_grad_loglik_vjp_wave
gpurec/fit/map_cv.py:25               from gpurec.api._execution     import stream_batches
\end{lstlisting}

Two of these import a private function inside a private module, from another
package.  \texttt{stream\_batches} is used by three packages, so it is shared
infrastructure living in a file marked private.  Either the underscore comes off
and the file moves somewhere shared, or the convention is dropped; as it stands
it tells the reader something untrue.

\subsection{\texorpdfstring{\dir{fit/} and \dir{solver/} overlap}{fit/ and solver/ overlap}}

The stated split --- \dir{solver/} is reusable mathematics, \dir{fit/} is
per-mode recipes --- reads well but is not held:

\begin{itemize}
  \item A mode-specific recipe lives in \dir{solver/}:
        \dir{solver/curvature/genewise.py}, 456 lines named after a fitting
        mode, alongside \dir{fit/genewise\_fit.py}.
  \item Newton drivers live in both, and \loc{fit/newton\_cg.py:82} imports one
        from \dir{solver/curvature/receiver.py}, so a fitting driver calls a
        solver driver.
  \item Generic machinery lives in \dir{fit/}: \dir{fit/optimize.py} describes
        itself as ``dataset-agnostic''.
\end{itemize}

\dir{fit/dtl\_fit.py} is the exception and the model to follow --- its docstring
claims to be ``the ONLY place that choice is made'', which is the right idea.

\dir{api/} versus \dir{core/} is defensible: 6 files of model wiring against 21
files of kernels and scheduling.  The only blur was the two imports resolved in
Section~\ref{sec:imports}.

\subsection{The two loose files at the package root}

\dir{gpurec/optimization.py} (96 lines) is load-bearing but homeless: it is
imported by \loc{gpurec/\_\_init\_\_.py:4-6} and two fitting modules, and it
pairs with \dir{gpurec/config/rates.py}, whose own docstring names it.  It
belongs beside that file.

\dir{gpurec/distributed.py} (112 lines) has zero importers anywhere and is
already counted in Section~\ref{sec:dead-code}.

\subsection{The public surface, and one packaging gap}

\dir{gpurec/\_\_init\_\_.py} is 7 lines exporting 7 names.  One of them,
\texttt{sample\_reconciliations}, is pulled straight from a \dir{core} internal
with no intermediate export point.

\dir{gpurec/config/defaults.toml} is \textbf{not shipped in the wheel} --- there
is no packaging entry for it, confirmed by building one.  The practical impact
is small, since that file is never read at runtime anyway, but it means the
drift check described at \loc{gpurec/config/gpurec\_config.py:200-202} cannot be
run against an installed copy.

\end{document}
